% !TEX root = ../paper.tex

\section{Experimental Setup}

\subsection{Baselines}

We evaluate both classical and learning-based baselines under an explicit off-the-shelf setting.
All learning-based methods use public pretrained weights or public pretrained model interfaces and
are evaluated without training or adaptation on UAV-TAlign-12K. This choice keeps the comparison
focused on practical deployability rather than dataset-specific adaptation.
Unified inference wrappers standardize model loading, image I/O, homography fitting, and output
schemas while preserving the public pretrained baseline parameters and settings.

The baseline suite includes a domain-specific infrared-visible classical method and modern
dense/learned matchers. RIFT2 is evaluated through a Python implementation under a resize-aware
setting. The modern pairwise baselines are Kornia LoFTR (pretrained="outdoor"), RoMa outdoor full,
XoFTR-640, and raw MINIMA. Raw MINIMA uses the MINIMA RoMa branch as a direct pairwise baseline;
reliability-guided evidence selection, robust scene consensus, and QA-aware verification are
activated in the UAV-TAlign full framework
\cite{li2020rift,sun2021loftr,edstedt2024roma,tuzcuoglu2024xoftr,ren2025minima}.

\subsection{Shared Geometric Fitting}

For the pairwise methods, the common downstream step is homography fitting from the predicted
correspondences. After each matcher returns point pairs and optional confidence values, the
shared evaluator estimates a homography with a robust backend and then records a unified outcome
category. This unified post-processing is important because it makes the pairwise baseline
table comparable at the level of usable registration outcomes rather than raw matcher outputs
alone.

The common evaluator uses a robust homography backend across the modern pairwise baselines:
\begin{itemize}
\item robust fitting method: USAC-MAGSAC;
\item reprojection threshold: 3.0;
\item confidence: 0.999;
\item maximum iterations: 10000;
\item spatial coverage grid: $4 \times 4$.
\end{itemize}

RIFT2 is reported as \emph{RIFT2 (Python, resize-aware)}. Its disclosed configuration is
\texttt{max\_dim=1200}, \texttt{npt=2000}, \texttt{patch\_size=64}, Lowe ratio 0.95, and OpenCV
USAC-MAGSAC for downstream homography estimation. We report it as a fixed Python resize-aware
deployment setting for the UAV-TAlign-12K evaluation.

\subsection{Main-Table Baseline Choices}

The main-table baseline choices are as follows. Raw MINIMA uses the MINIMA RoMa backend
with the \texttt{minima\_roma} checkpoint family and the large/full RoMa branch. RoMa uses the
public outdoor full model. XoFTR uses the 640-resolution visible-thermal checkpoint family as the
default main-table setting. Kornia LoFTR (pretrained="outdoor") uses the public pretrained outdoor
model from Kornia.

\subsection{Evaluation Details}

The LoFTR baseline is evaluated under a resize-aware inference setting. The longest image
dimension used for matching is limited to 1600, and automatic mixed precision is disabled. This
deployment setting preserves the public pretrained model while supporting consistent full-subset
evaluation within the fixed deployment budget.

All experiments were run in a Linux Conda environment with Python 3.10,
PyTorch 2.0.1, torchvision 0.15.2, and CUDA 11.8 for the main learned baselines and UAV-TAlign
pipeline. RIFT2 is executed in an independent Python environment. All evaluation outputs are
written to isolated result roots, with raw dataset directories and source trees treated as
read-only inputs during evaluation.
